\section{A linear bound for stabilization in rank-two graph towers}
\label{sec:rank-two-onset}
Let $p$ be prime. A rank-two integral voltage tower attaches a vector of
$\mathbb Z^2$ to each oriented base edge, with negation on reversal. At
level $n$ the deck group is $(\mathbb Z/p^n\mathbb Z)^2$ and following
an edge adds its voltage to the deck coordinate. Write $v_n$ for the
$p$-adic valuation of the number of spanning trees. DuBose and
Vallières proved eventual polynomiality of $v_n$ in $p^n,n$
(\href{https://doi.org/10.5802/alco.304}{Theorem~6.2}).
The preceding \href{https://github.com/VladimirReshetnikov/ProveIt/blob/bd8864aa31430944f587f544213d562665f1d817/docs/incoming/Effective_Graph_Towers.zip}{ProveIt graph-tower report}
gives an algorithm in every rank and asks, in its first further research
question, for the growth of the worst possible onset as a function of
voltage bit length. We obtain a sharp order in rank two at $p=3$ for a
bouquet of five loops. The upper estimate works for all primes and
rank-two Laurent determinants. It uses the cyclotomic identity
$v_p(\zeta_{p^h}-1)=1/\varphi(p^h)$, also used in
\href{https://github.com/lsernau/grh-lean/blob/72fead102f663d5c7a407964a13dc0cf9d977207/grh/src/New/Congruences/Selmer/RootsOfUnity.lean}{the repository's specialization argument}.
The new ingredient is separation of the finitely many rational slopes
where pairs of Laurent monomials coincide.

\paragraph{Quantitative shell theorem.}
Let $F\ne0$ be an integral Laurent polynomial in two variables,
supported in $[-H,H]^2$, with $H\ge1$. Assume $F$ is nonzero at all
nontrivial pairs of $p$-power roots of unity. For $h\ge1$ write
$S_h$ for the sum of $v_p(F(x,y))$ over pairs of joint order exactly
$p^h$. Set
\[
 K=\lceil\log_p(8H^2+1)\rceil,\quad
 M=\lceil\log_p(2H+1)\rceil,\quad h_0=K+M+2.
\]
There exist rational $s_2,s_1,t_1,s_0$ such that for every $h\ge h_0$,
\[
 S_h=s_2p^{2h}+(s_1+t_1 h)p^h+s_0.
\]
All constants are independent of the level. The bound $h_0$ only
depends on the exponents and $p$.

\paragraph{Proof.}
Normalize $v_p(p)=1$ and write $e_h=(p-1)p^{h-1}$.
Dividing $F$ by the common $p$-power $p^g$ of its coefficients adds
$g(p^{2h}-p^{2h-2})$ back to the shell. Collect equal exponents and
assume henceforth that $F=\sum_i c_i X^{a_i}Y^{b_i}$ has a unit
coefficient. Fix compatible primitive roots $\zeta_h$. Simultaneous
raising to units modulo $p^h$ preserves the valuation, with each
orbit of size $e_h$. Its contribution is therefore
$\operatorname{ord}_h(F)=e_h v_p(F)$. Each orbit is uniquely
represented in one of two charts, $(\zeta_h,\zeta_h^z)$ for any $z$
modulo $p^h$, or $(\zeta_h^z,\zeta_h)$ with $p\mid z$. In the
second chart interchange the two exponent coordinates below.

The projected exponents in the first chart are $a_i+b_i z$. Take as
centers the distinct $p$-integral rational solutions of
$a_i+b_i z=a_j+b_j z$ lying in the chart. Write each center as a
reduced fraction $A/B$, with $p\nmid B$, $|A|,|B|\le2H$. The
difference of two distinct centers has nonzero numerator of absolute
value at most $8H^2$. Hence their disks $z\equiv A/B\pmod{p^K}$
are disjoint. In the restricted chart each disk is entirely inside
or outside the chart since $K\ge1$.

Outside these disks the projected exponents are pairwise distinct
modulo $p^J$, where $J=K+M$. In fact a nonzero slope difference
$b_i-b_j$ has valuation $<M$, and the distance factor
$z-(a_j-a_i)/(b_i-b_j)$ has valuation $<K$; this covers centers
outside the chart too. A nonzero constant difference has valuation
$<M$. For each fixed $z\bmod p^J$ in this complement let
$n_i\in[0,p^J)$ be the distinct exponent residues. The polynomial
$\sum_i c_i(1+T)^{n_i}$ is nonzero modulo $p$ by the distinct
degrees and the unit coefficient. Its first nonzero coefficient
occurs at $j<p^J$. For any lift of $z$ and $h\ge h_0$, the orbit
contributes exactly $j$. Indeed, modulo $(p,T^{p^J})$ the binomial
expansion of each monomial depends only on its exponent modulo $p^J$.
One may choose nonnegative representatives modulo $p^h$ in forming
the actual polynomial. At $T=\zeta_h-1$ a coefficient divisible
by $p$ has order at least $e_h\ge p^J$, and all later terms have
larger order than $j$. Each class contributes $j p^{h-J}$ to $S_h$.

Fix a center $A/B$. Applying the Galois automorphism
$\zeta_h\mapsto\zeta_h^B$ expresses the same valuation as the one of
\[
 P(\zeta_h,W),\qquad
 W=\zeta_h^{Bz-A},\qquad
 P(X,Y)=\sum_i c_i X^{B a_i+A b_i}Y^{b_i},
\]
up to multiplication by a monomial of valuation zero. Shift each
exponent axis by its minimum so that $P$ is a polynomial of degrees
$D\le8H^2$ in $X$ and $E\le2H$ in $Y$. The exponent transformation
is injective, so $P$ still has a unit coefficient. As $z$ ranges over
the disk modulo $p^h$, $W$ ranges exactly once over the roots of
order dividing $p^{h-K}$, since $B$ is a unit. Write $r$ for the
conductor exponent of $W$, with $r=0$ for $W=1$, and put $R=M+1$.

For $0\le r\le R$, use $e_0=1$ and write
$P(1+T,W)=\sum_{j=0}^{D} d_j(W)T^j$.
This polynomial is nonzero: its vanishing would give a torsion zero
of $F$ at the pair represented by $z$. Its nonzero coefficients have
valuations in $e_r^{-1}\mathbb Z$. Let $\mu_W$ be their minimum
and let $j_W$ be the first index attaining it. We have
$e_h/e_r>D$ (for $r\ge1$ the quotient is $p^{h-r}$ and for $r=0$
use $e_h\ge p^{h-1}$). Thus terms with larger coefficient valuation
cannot minimize. The unique minimum has order $e_h\mu_W+j_W$.
The set of all $W$ with $r\le R$ is fixed for every $h\ge h_0$.
Its total is a constant multiple of $p^h$ plus a constant.

For $R<r\le h-K$, put $k=h-r$. Then
$\operatorname{ord}_h(W-1)=p^k>D$ and $e_r>E+1$.
Expand
$P(X,1+U)=\sum_{i=0}^E P_i(X)U^i$, and take the smallest $l$
such that $P_l(X)$ is nonzero modulo $p$. Such an $l$ exists.
Let $d\le D$ be the order of $P_l(1+T)$ at $T=0$ modulo $p$.
Its value at $X=\zeta_h$ has order exactly $d$. Terms with $i<l$
have order at least $e_h=e_rp^k>lp^k+d$. Terms with $i>l$ have
order at least $(l+1)p^k>lp^k+d$. The orbit contribution is exactly
$lp^k+d$ for every root of conductor $p^r$. The contribution from
all these roots is
\[
 \sum_{k=K}^{h-R-1}e_{h-k}(lp^k+d)
 =l(h-R-K)e_h+d(p^{h-K}-p^R).
\]
Adding the finitely many disks in each chart, their complements,
and the removed content proves the shell formula.

\paragraph{Consequence for graph towers.}
Let $F$ be the determinant of the Laurent voltage Laplacian of a
finite connected base graph, and assume the finite lifts are
connected. Diagonalizing by the characters of the deck group splits
the lifted Laplacian into blocks. The trivial character yields the
base Laplacian; connectedness gives invertible blocks at every
nontrivial character, since the lifted Laplacian has nullity one.
Applying the matrix-tree identity to the base and the lift gives
\[
 v_n=v_0-2n+\sum_{h=1}^{n} S_h .
\]
Summing the shell formula shows that from $n=h_0$ onward
\[
 v_n=A p^{2n}+(B+C n)p^n+D' n+E'
\]
for rational constants. For a bouquet the support of its Laurent
Laplacian consists of zero and the signed voltage vectors. Thus if
the absolute values of the coordinates are at most $2^b$, the
first stable level is at most $3b\log_p 2+O_p(1)$. For a fixed base
graph the coordinates in the determinant support are bounded by
the vertex count times the largest voltage coordinate, by determinant
expansion. The same bit-length estimate follows with a base-dependent
additive constant. Five consecutive values beyond this upper bound
suffice to recover the law: the evaluation matrix on
$1,n,p^n,np^n,p^{2n}$ is nonsingular since these sequences solve the
order-five recurrence with characteristic polynomial
$(x-1)^2(x-p)^2(x-p^2)$ and are linearly independent. Checking the
finitely many preceding levels then gives the least onset. This is a
bound on conductor precision; enumerating all the characters there
can still be expensive.

\paragraph{Matching lower bound at $p=3$.}
The preceding graph report gives five-loop bouquets with directions
$3^s e_1,e_2,e_1,e_1,e_1$ for $s\ge2$.
For completeness let $A(U)=2-U-U^{-1}$. At joint conductor $3^j$
the sum $A(U)+A(V)$ has valuation exactly $2/e_j$: its normalized
quadratic coefficient in $\zeta_j-1$ is $-(a^2+b^2)$ for exponents
$a,b$ not both divisible by $3$, which is nonzero modulo $3$.
Consequently $A(X^{3^s})+A(Y)$ dominates $3A(X)$ off the
kernel of $(X,Y)\mapsto(X^{3^s},Y)$, since its valuation is at most
one. Put $t=\min(n,s)$. The image is $C_{3^{n-t}}\times C_{3^n}$
and each fiber has size $3^t$. Its order-$3^j$ shells have sizes
$3^{2j}-3^{2j-2}$ for $j\le n-t$ and $3^{n-t}e_j$ for $n-t<j\le n$.
The kernel contributes $3^t-1+2t$, using
$\prod_{X^{3^t}=1,\ X\ne1}(1-X)=3^t$.
Adding these contributions and subtracting $2n$ gives
\[
 v_n=2t3^n+4(3^n-3^t)+(3^t-1)+2t-2n.
\]
At $n\ge s$ this is the eventual polynomial obtained by substituting
$t=s$. At $n=s-1$ its actual value exceeds the extrapolated polynomial
by $4\cdot3^{s-1}-2$. The least onset is therefore $s$. With
$3^s\le2^b$, the worst-case onset for this fixed prime and loop count
grows as $\Theta(b)$. The sharp leading constant and higher ranks
remain questions.

\paragraph{Sharp order at every prime for fixed loop count.}
The coefficient of $p^{2n}$ in the graph formula is exactly the
common $p$-adic content $g=\min_i v_p(c_i)$ of its Laurent
determinant, in the sense of the valuation of the content. Indeed,
restoring the removed power contributes exactly $g(p^{2n}-1)$
to the cumulative sum. Every other summand grows as $O(np^n)$.
This yields a second lower bound that uses no restriction on $p$.
On a bouquet of $2p+1$ loops put $p$ copies each of $e_1,e_2$ and
one step $p^s e_1$, with $s\ge1$. Compare with the same assignment
where the final step is zero. Both towers are connected, and the
graphs agree through level $s$. The zero-step determinant is
$p(A(X)+A(Y))$ and has content valuation one. The extra summand
$A(X^{p^s})$ changes the content valuation to zero because its
two nonconstant exponents are new. The two eventual laws have
different $p^{2n}$ coefficients. For $s$ large enough that
$s-4\ge h_0(1)$, the zero-step law is already exact throughout
$[s-4,s]$. If the other law were exact from $s-4$ onward the
two would agree on these five consecutive levels, which the
recurrence argument excludes. Its least onset is at least $s-3$.
Taking $s=\lfloor\log_p(2^b)\rfloor$ gives worst-case onset
$\Theta_p(b)$ on this one fixed bouquet for every prime. Padding
with zero-voltage loops gives the same order for any larger fixed
loop count. The five-loop ternary calculation above gives a
smaller loop count at that prime and identifies the exact onset.

\paragraph{Classification by number of loops.}
In fact the threshold loop count can be found exactly. Let $L_{p,m}(b)$
be the largest least onset over connected rank-two towers on one bouquet
of $m\ge2$ loops, with coordinates bounded by $2^b$ in absolute value.
Repeated directions and zero voltages are allowed. As $b\to\infty$,
\[
 \boxed{\quad L_{p,m}(b)=
 \begin{cases}
 O_p(1),&2\le m<\max(3,p),\\
 \Theta_p(b),&m\ge\max(3,p).
 \end{cases}\quad}
\]
Constants here can also be taken independent of $m$ in the upper
bound. For odd $p$ and $m<p$, the onset is at most one.

To see the bounded range for odd $p$, take a projective direction
modulo $p$. The projected step residues are not all zero because
connectedness forces the steps to span $\mathbb F_p^2$. In the group
ring $\mathbb F_p[Z]/(Z^p-1)$ take any nonzero residue $a$ occupied
by a step, and let $j$ count the steps in $\{a,-a\}$. The
coefficient of $Z^a$ in the projected Laplacian is $-j\ne0$
since $a\ne-a$ and $1\le j\le m<p$. Thus its expansion at $Z=1$
has first surviving order $<p$. Since $e_h\ge p$ for $h\ge2$,
every orbit of depth $h$ lifting that projective direction has exactly
that integral order. Every residue direction has $p^{h-1}$ lifts,
so $S_h$ is a fixed multiple of $p^{h-1}$ for $h\ge2$. Summation
from level one proves the claim.
For $p=2,m=2$ the two steps form an invertible basis modulo $2$,
hence modulo $2^n$ for every $n$. A change of coordinates gives the
same square grid at every fixed level for all choices of voltages,
so the support bound $h_0$ evaluated at $H=1$ suffices.

For the lower range it is enough to work at the threshold value
and pad with zero voltages. We use a test supplied by the shell proof.
If a Laurent polynomial is primitive modulo $p$, the coefficient of
$np^n$ in the cumulative formula is positive precisely when some
nonzero rational direction $(u,v)$ satisfies
\[
                F(Z^u,Z^v)\equiv0\pmod p
                \quad\text{as a Laurent polynomial in }Z,
\]
where $(u,v)$ is taken to be a primitive integer representative.
Indeed the coefficient of $h p^h$ in $S_h$ is the sum of the
nonnegative disk contributions $l(p-1)/p$. A disk contributes
positively exactly when $P(X,1)$ vanishes modulo $p$. This is the
displayed identity after the monomial shift and the change of
coordinates defining the disk. Conversely any such identity must
group at least two distinct exponents with the same projection
because there is a unit coefficient. This yields one of the
rational centers, using the chart where the direction has a unit
pivot. Summing $hp^h$ gives a positive multiple of $np^n$
plus a multiple of $p^n$ and a constant. This proves the test.

Suppose $m=p$ is odd. Use $p-2$ copies of $(1,0)$, one copy
of $(1,1)$ and one copy of $(1+\delta,2)$, first with $\delta=0$
and then with $\delta=p^s>4$. In both cases the content is a unit
and the first two directions ensure connectedness. In the first
case projection along $(u,v)=(1,0)$ makes all $p$ exponents equal
to one, so the test is positive. For the second case a vanishing
projection must have all $p$ absolute projected exponents equal
to one another and nonzero. Indeed an occupied positive absolute
frequency $d$ contributes the negative of its multiplicity to the
coefficient of $Z^d$. This multiplicity must equal $p$, and a
nonzero direction cannot annihilate every step. Thus $u\ne0$ and
$|u+v|=|u|$, so $v=0$ or $v=-2u$. The last projected absolute
value is now $|(1+\delta)u|$ or $|(\delta-3)u|$, both different
from $|u|$. The $np^n$ term vanishes after perturbation.

For $p=2,m=3$ use the three steps
$(1,0),(1,1),(\delta,1+2\delta)$ with $\delta=0$ or $2^s$.
Again the content is a unit. In characteristic two every positive
absolute projected frequency must occur an even number of times.
Three steps spanning the plane therefore give a zero Laurent
polynomial exactly when one projected value is zero and the other
two have equal nonzero absolute value. A direction annihilating
one step projects the other two in proportion to their determinants
with that step. For $\delta=0$ the three pairwise absolute
determinants all equal one, and such a direction exists. For $\delta>0$
the determinants are $1,1+\delta,1+2\delta$, all distinct.
Thus this pair also has different $np^n$ coefficients.

In both constructions the perturbed tower agrees exactly with the
fixed initial tower up through level $s$. Once the fixed initial
tower has stabilized by $s-4$, the five-consecutive-values argument
forces onset at least $s-3$ for the perturbed tower. The heights
are at most $1+2p^s$, allowing $s=b\log_p 2+O_p(1)$. Together with
the upper bound this proves the classification. It leaves open the
optimal leading constant within the linear-growth range.
